\chapter*{Lời Tựa}
\addcontentsline{toc}{chapter}{Lời Tựa}
\markboth{Lời Tựa}{Lời Tựa}

\begin{truyen}{Lời Tựa}{Witty Classroom Replies}
\chuhoa{C}{uốn này đi theo một kiểu rất đời lớp học: học sinh hỏi những câu quá quen, còn giáo viên đáp bằng kiểu cà khịa nhẹ, thâm một chút, nhưng vẫn đủ duyên để người đọc cười trước rồi ngẫm sau.}
\textit{(This book follows a very familiar classroom rhythm: students ask predictable questions, and teachers respond with light sarcasm, subtle wit, and a smile that hides a lesson.)}

Điều quan trọng của cuốn này không nằm ở chữ ``chửi''. Nó nằm ở chữ ``duyên''. Không phải chửi nặng, không hạ nhục, không cay nghiệt. Chỉ là những câu đáp vừa đủ sắc để học sinh thấy mình bị bóc, nhưng không đến mức bị tổn thương.
\textit{(The key word here is not insult but wit. The replies are not cruel or humiliating. They are just sharp enough to expose a weak excuse without wounding the learner.)}

Trong 10 chương và 1000 màn hỏi -- đáp, người đọc sẽ gặp đủ kiểu câu rất quen: chưa hiểu, quên làm bài, xin điểm, sợ lên bảng, hỏi bài này để làm gì, hỏi có thi không, hỏi vì sao bị phạt, hỏi tương lai rồi sẽ ra sao. Và phía sau mỗi câu cà khịa của thầy, luôn có một bài học nhỏ về kỷ luật, tư duy, trách nhiệm và sự trưởng thành.
\textit{(Across 10 chapters and 1000 exchanges, readers will meet the most familiar school questions: not understanding, forgetting homework, asking for points, fearing the board, asking what a lesson is for, asking whether it will be on the test, questioning discipline, and worrying about the future. Behind each witty teacher reply sits a small lesson about responsibility and maturity.)}

Nếu bạn là học sinh, mong bạn đọc để thấy mình đâu đó trong những câu hỏi này và cười được vào chính mình. Nếu bạn là giáo viên, mong bạn đọc để có thêm chất liệu rất gần gũi khi đối thoại với học trò.
\textit{(If you are a student, may you see yourself somewhere in these questions and learn to laugh at your own excuses. If you are a teacher, may this book offer useful language for speaking with learners.)}
\end{truyen}

\begin{baihoc}
Cà khịa đúng cách không làm lớp học lạnh đi. Nó làm bài học đỡ trôi tuột.
\textit{(Wit, used well, does not make a classroom colder. It makes the lesson harder to forget.)}
\end{baihoc}

\ngancach